\begin{table}[htb!]
  \centering
  \tabcaption{Effect of Flash Floods on Hazard of Establishment Closure}
  \label{tab6: hazard_closure}
  \scalebox{0.75}{
 \begin{threeparttable}
    \begin{tabular}{lcccc}
    \toprule
          & \multicolumn{2}{c}{(1)} & \multicolumn{2}{c}{(2)} \\
          & Coefficient & Hazard Ratio & Coefficient & Hazard Ratio \\
    \midrule
    Flash Flood Post & 0.06067*** & 1.06255*** & 0.05868*** & 1.06044*** \\
                      & (0.01360) &  & (0.01324) &  \\
    Observations & 367{,}402 & 367{,}402 & 367{,}402 & 367{,}402 \\
    Number of Units & 79{,}604 & 79{,}604 & 79{,}604 & 79{,}604 \\
    Number of Failures & 39{,}340 & 39{,}340 & 39{,}340 & 39{,}340 \\
    Establishment Controls & No & No & Yes & Yes \\
    \bottomrule
    \end{tabular}%
    \begin{tablenotes}[flushleft]
    \item \small Notes: Exact R translation of the legacy Stata specification used for Table 6. The survival setup follows \texttt{stset year, failure(morte) id(id\_estab)} and the Cox model follows \texttt{stcox} with sector-year strata and establishment-clustered standard errors. Column (2) adds the baseline controls generated in the Stata do-file from each establishment's first observed year: employment (\texttt{empregados2}), branch count (\texttt{afil2}), opening year (\texttt{yr\_abert2}), and workforce education shares (\texttt{educ\_d22}, \texttt{educ\_d32}, \texttt{educ\_d42}). Observations denote establishment-year observations used by each model, the number of units corresponds to unique establishments at risk, and the number of failures indicates establishments that experienced closure. *** p$<$0.01, ** p$<$0.05, * p$<$0.1.
    \end{tablenotes}
    \end{threeparttable}
    }
\end{table}%
